\section{Agent：短程高频反馈与 LM 内核}

星眠与 Dokie 讲完后，重点自然落到 Dokie 背后的 Agent 判断。本节把张月光的 Agent know-how 拆成产品路径和技术内核两部分。

访谈后半段进入 Agent。张月光给出两个 know-how。第一，不要只把 Agent 当成长程离线任务，越长程越不稳定；也许应该做更短程、更高频、更贴近用户反馈的循环。第二，至少今天，LM 仍然是 Agent 的智力内核，因为它能理解用户意图、拆解任务、调用工具和调模型。

\begin{figure}[H]
\centering
\includegraphics[width=0.96\textwidth]{agent-short-loop.png}
\caption{Agent 产品的短程高频反馈路线。}
\end{figure}

\begin{knowledgebox}{读图：短程高频不是能力低，而是产品可控}
图中用户意图进入 LM 内核，LM 调工具和模型生成产物，用户再给高频反馈。相比把 Agent 放进长程任务里等待最终结果，短程高频更容易建立信任、降低延迟焦虑，并把用户判断纳入闭环。
\end{knowledgebox}

\subsection{LM 为什么仍是内核}

他不太认同“语言无法达到智能终点”的说法，甚至认为语言就是智能的本质。这个观点和 EP133 谢赛宁的世界模型观点形成有趣对照。张月光站在产品和 Agent 视角：用户意图主要以语言表达，工具调用需要语言规划，工作产物也大量是语言或结构化内容，因此 LM 是当前最自然的治理内核。

\begin{warningbox}{这里有一组有价值的分歧}
谢赛宁强调语言不能替代物理世界；张月光强调语言是 Agent 的智能内核。二者并不必然冲突：物理世界智能需要超越语言，生产力 Agent 现阶段则高度依赖语言组织任务。
\end{warningbox}

\subsection{本章小结}

Agent 产品不是只拼长程 benchmark。短程、高频、低延迟、用户反馈和可编辑结果，可能更早形成真实产品体验。LM 仍是当下 Agent 的核心大脑。

